\documentclass[11pt]{article}
\usepackage[a4paper,margin=1in]{geometry}
\usepackage{subfiles}
\usepackage{amsmath,amssymb,bm}
\usepackage{hyperref}
\usepackage{booktabs}
\usepackage{array}
\usepackage{mathtools}

\title{Appendix 1 (to main theory): Field Equations from the One--Metric ISPG Action}
\author{}
\date{}

\begin{document}
\let\phi\varphi
\maketitle

% Reset section counter for inclusion in master document
\setcounter{section}{0}

\section*{Purpose and scope}
This appendix derives the vacuum field equations of the main theory from the covariant one-metric action, and records the sourced scalar equation under minimal coupling. The output is a compact, citable set of equations used throughout the main theory appendix package.

\section*{Non-negotiable consistency with the main theory (one-metric formulation)}
We work strictly within the main theory posture:
\begin{itemize}
\item matter is minimally coupled to the single physical metric $g_{\mu\nu}$ (no disformal mapping and no frame change);
\item the covariant backbone is
\begin{equation}
S=\frac{1}{16\pi G}\int d^4x\sqrt{-g}\left[R+\frac12\,g^{\mu\nu}\partial_\mu\phi\,\partial_\nu\phi\right]+S_m[g_{\mu\nu},\psi],
\label{eq:action_app1}
\end{equation}
with the same metric $g_{\mu\nu}$ as used everywhere else in the main theory.
\end{itemize}

\section{Variational identities}\label{sec:idents}
We use standard variations:
\begin{align}
\delta\sqrt{-g}&=-\frac12\sqrt{-g}\,g_{\mu\nu}\,\delta g^{\mu\nu},\\
\delta(\sqrt{-g}R)&=\sqrt{-g}\left(G_{\mu\nu}\,\delta g^{\mu\nu}+\nabla_\alpha \Theta^\alpha\right),
\end{align}
where $\Theta^\alpha$ is a total-derivative boundary term that vanishes for suitable boundary conditions.
For the scalar kinetic term,
\begin{equation}
\delta\!\left[\sqrt{-g}\,\frac12 g^{\mu\nu}\partial_\mu\phi\,\partial_\nu\phi\right]
=\sqrt{-g}\left[\frac12\partial_\mu\phi\,\partial_\nu\phi-\frac14 g_{\mu\nu}(\partial\phi)^2\right]\delta g^{\mu\nu}
-\sqrt{-g}\,\nabla_\mu(\partial^\mu\phi)\,\delta\phi,
\label{eq:scalarvar}
\end{equation}
after integrating by parts in the $\delta\phi$ variation (the boundary term vanishes for suitable boundary conditions).

\section{Metric field equations}\label{sec:metric_app1}
Define the matter stress-energy tensor by
\begin{equation}
\delta S_m=-\frac12\int d^4x\sqrt{-g}\,T^{(m)}_{\mu\nu}\,\delta g^{\mu\nu}.
\label{eq:Tmdef_app1}
\end{equation}
Varying \eqref{eq:action_app1} with respect to $g^{\mu\nu}$ yields
\begin{equation}
G_{\mu\nu}+S_{\mu\nu}=8\pi G\,T^{(m)}_{\mu\nu},
\label{eq:einsteinS_app1}
\end{equation}
with
\begin{equation}
S_{\mu\nu}\equiv \frac12\left(\partial_\mu\phi\,\partial_\nu\phi-\frac12 g_{\mu\nu}(\partial\phi)^2\right).
\label{eq:Smunu_app1}
\end{equation}
Equivalently, introducing the canonical-form scalar tensor
\begin{equation}
T^{(\phi)}_{\mu\nu}\equiv \partial_\mu\phi\,\partial_\nu\phi-\frac12 g_{\mu\nu}(\partial\phi)^2,
\label{eq:Tphi_app1}
\end{equation}
The main theory bookkeeping is
\begin{equation}
G_{\mu\nu}=8\pi G\,T^{(m)}_{\mu\nu}+\kappa\,T^{(\phi)}_{\mu\nu},
\qquad \kappa=-\frac12,
\label{eq:einsteinT_app1}
\end{equation}
so that $S_{\mu\nu}=-\kappa T^{(\phi)}_{\mu\nu}$.

\section{Scalar field equation}\label{sec:scalar}
Varying \eqref{eq:action_app1} with respect to $\phi$ gives the vacuum equation
\begin{equation}
\square\phi=0 \qquad (\text{vacuum}).
\label{eq:boxphi_app1}
\end{equation}
Under minimal coupling ($\delta S_m/\delta\phi\equiv 0$), direct $\phi$-variation alone yields \eqref{eq:boxphi_app1} also in the presence of matter. The sourced form used throughout the main theory,
\begin{equation}
\square\phi=-\frac{8\pi G}{c^4}\,T,
\qquad T\equiv g^{\mu\nu}T^{(m)}_{\mu\nu},
\label{eq:boxphiT}
\end{equation}
is therefore \emph{not} an independent Euler--Lagrange equation; it is the reduced-sector consistency relation obtained by tracing the metric equation \eqref{eq:einsteinT_app1} after substituting the bi--conformal ansatz $g_{\mu\nu}[\phi]$ (Appendix~3). In matter-free regions ($T=0$) it coincides with \eqref{eq:boxphi_app1}, so the weak-field exterior used in PPN/Shapiro/lensing analyses is unaffected by the distinction.

\section{Conservation identities (consistency loop)}\label{sec:cons_app1}
Taking the divergence of \eqref{eq:einsteinT_app1} and using $\nabla^\mu G_{\mu\nu}=0$ gives
\begin{equation}
\nabla^\mu T^{(m)}_{\mu\nu} = -\frac{\kappa}{8\pi G}\,\nabla^\mu T^{(\phi)}_{\mu\nu}.
\label{eq:balance_app1}
\end{equation}
Minimal coupling implies $\nabla^\mu T^{(m)}_{\mu\nu}=0$, and the scalar identity
\begin{equation}
\nabla^\mu T^{(\phi)}_{\mu\nu}=(\square\phi)\,\partial_\nu\phi
\label{eq:divTphi_app1}
\end{equation}
closes the system on-shell.

\section{Derived vs.\ stated: compact ledger}\label{sec:ledger_app1}
\begin{center}
\begin{tabular}{@{}p{0.28\linewidth}p{0.62\linewidth}@{}}
\toprule
\textbf{Derived here} &
Metric field equation \eqref{eq:einsteinS_app1}--\eqref{eq:einsteinT_app1}; explicit scalar contribution \eqref{eq:Smunu_app1}--\eqref{eq:Tphi_app1}; vacuum scalar equation \eqref{eq:boxphi_app1}; conservation loop \eqref{eq:balance_app1}--\eqref{eq:divTphi_app1}.\\
\midrule
\textbf{Imposed (reduced-sector consistency)} &
Sourced form \eqref{eq:boxphiT}: not an independent Euler--Lagrange equation, but the trace consistency relation of \eqref{eq:einsteinT_app1} after the bi--conformal ansatz of Appendix~3.\\
\midrule
\textbf{Used elsewhere} &
PPN/deflection/perihelion (Appendices 4,5,7) use the weak-field exterior of \eqref{eq:einsteinT_app1}; Shapiro 2PN uses the bi--conformal exterior solution; GW and stability appendices use minimal coupling and principal-symbol structure; Appendix 2 provides the detailed $T^{(\phi)}_{\mu\nu}$ discussion.\\
\bottomrule
\end{tabular}
\end{center}

\section*{Takeaway}
From the main theory one-metric action \eqref{eq:action_app1}, direct variation yields the metric equation
\[
G_{\mu\nu}=8\pi G\,T^{(m)}_{\mu\nu}-\frac12\,T^{(\phi)}_{\mu\nu},
\]
and the vacuum scalar equation $\square\phi=0$. The trace-consistency form $\square\phi=-(8\pi G/c^4)T$ used in the main theory is not an independent Euler--Lagrange equation; it is the reduced-sector consistency relation obtained after substituting the bi--conformal ansatz $g_{\mu\nu}[\phi]$ (Appendix~3). Covariant conservation closes on-shell under minimal coupling.

\end{document}
